% Options for packages loaded elsewhere
% Options for packages loaded elsewhere
\PassOptionsToPackage{unicode}{hyperref}
\PassOptionsToPackage{hyphens}{url}
\PassOptionsToPackage{dvipsnames,svgnames,x11names}{xcolor}
%
\documentclass[
  letterpaper,
  DIV=11,
  numbers=noendperiod]{scrartcl}
\usepackage{xcolor}
\usepackage{amsmath,amssymb}
\setcounter{secnumdepth}{-\maxdimen} % remove section numbering
\usepackage{iftex}
\ifPDFTeX
  \usepackage[T1]{fontenc}
  \usepackage[utf8]{inputenc}
  \usepackage{textcomp} % provide euro and other symbols
\else % if luatex or xetex
  \usepackage{unicode-math} % this also loads fontspec
  \defaultfontfeatures{Scale=MatchLowercase}
  \defaultfontfeatures[\rmfamily]{Ligatures=TeX,Scale=1}
\fi
\usepackage{lmodern}
\ifPDFTeX\else
  % xetex/luatex font selection
\fi
% Use upquote if available, for straight quotes in verbatim environments
\IfFileExists{upquote.sty}{\usepackage{upquote}}{}
\IfFileExists{microtype.sty}{% use microtype if available
  \usepackage[]{microtype}
  \UseMicrotypeSet[protrusion]{basicmath} % disable protrusion for tt fonts
}{}
\makeatletter
\@ifundefined{KOMAClassName}{% if non-KOMA class
  \IfFileExists{parskip.sty}{%
    \usepackage{parskip}
  }{% else
    \setlength{\parindent}{0pt}
    \setlength{\parskip}{6pt plus 2pt minus 1pt}}
}{% if KOMA class
  \KOMAoptions{parskip=half}}
\makeatother
% Make \paragraph and \subparagraph free-standing
\makeatletter
\ifx\paragraph\undefined\else
  \let\oldparagraph\paragraph
  \renewcommand{\paragraph}{
    \@ifstar
      \xxxParagraphStar
      \xxxParagraphNoStar
  }
  \newcommand{\xxxParagraphStar}[1]{\oldparagraph*{#1}\mbox{}}
  \newcommand{\xxxParagraphNoStar}[1]{\oldparagraph{#1}\mbox{}}
\fi
\ifx\subparagraph\undefined\else
  \let\oldsubparagraph\subparagraph
  \renewcommand{\subparagraph}{
    \@ifstar
      \xxxSubParagraphStar
      \xxxSubParagraphNoStar
  }
  \newcommand{\xxxSubParagraphStar}[1]{\oldsubparagraph*{#1}\mbox{}}
  \newcommand{\xxxSubParagraphNoStar}[1]{\oldsubparagraph{#1}\mbox{}}
\fi
\makeatother

\usepackage{color}
\usepackage{fancyvrb}
\newcommand{\VerbBar}{|}
\newcommand{\VERB}{\Verb[commandchars=\\\{\}]}
\DefineVerbatimEnvironment{Highlighting}{Verbatim}{commandchars=\\\{\}}
% Add ',fontsize=\small' for more characters per line
\usepackage{framed}
\definecolor{shadecolor}{RGB}{241,243,245}
\newenvironment{Shaded}{\begin{snugshade}}{\end{snugshade}}
\newcommand{\AlertTok}[1]{\textcolor[rgb]{0.68,0.00,0.00}{#1}}
\newcommand{\AnnotationTok}[1]{\textcolor[rgb]{0.37,0.37,0.37}{#1}}
\newcommand{\AttributeTok}[1]{\textcolor[rgb]{0.40,0.45,0.13}{#1}}
\newcommand{\BaseNTok}[1]{\textcolor[rgb]{0.68,0.00,0.00}{#1}}
\newcommand{\BuiltInTok}[1]{\textcolor[rgb]{0.00,0.23,0.31}{#1}}
\newcommand{\CharTok}[1]{\textcolor[rgb]{0.13,0.47,0.30}{#1}}
\newcommand{\CommentTok}[1]{\textcolor[rgb]{0.37,0.37,0.37}{#1}}
\newcommand{\CommentVarTok}[1]{\textcolor[rgb]{0.37,0.37,0.37}{\textit{#1}}}
\newcommand{\ConstantTok}[1]{\textcolor[rgb]{0.56,0.35,0.01}{#1}}
\newcommand{\ControlFlowTok}[1]{\textcolor[rgb]{0.00,0.23,0.31}{\textbf{#1}}}
\newcommand{\DataTypeTok}[1]{\textcolor[rgb]{0.68,0.00,0.00}{#1}}
\newcommand{\DecValTok}[1]{\textcolor[rgb]{0.68,0.00,0.00}{#1}}
\newcommand{\DocumentationTok}[1]{\textcolor[rgb]{0.37,0.37,0.37}{\textit{#1}}}
\newcommand{\ErrorTok}[1]{\textcolor[rgb]{0.68,0.00,0.00}{#1}}
\newcommand{\ExtensionTok}[1]{\textcolor[rgb]{0.00,0.23,0.31}{#1}}
\newcommand{\FloatTok}[1]{\textcolor[rgb]{0.68,0.00,0.00}{#1}}
\newcommand{\FunctionTok}[1]{\textcolor[rgb]{0.28,0.35,0.67}{#1}}
\newcommand{\ImportTok}[1]{\textcolor[rgb]{0.00,0.46,0.62}{#1}}
\newcommand{\InformationTok}[1]{\textcolor[rgb]{0.37,0.37,0.37}{#1}}
\newcommand{\KeywordTok}[1]{\textcolor[rgb]{0.00,0.23,0.31}{\textbf{#1}}}
\newcommand{\NormalTok}[1]{\textcolor[rgb]{0.00,0.23,0.31}{#1}}
\newcommand{\OperatorTok}[1]{\textcolor[rgb]{0.37,0.37,0.37}{#1}}
\newcommand{\OtherTok}[1]{\textcolor[rgb]{0.00,0.23,0.31}{#1}}
\newcommand{\PreprocessorTok}[1]{\textcolor[rgb]{0.68,0.00,0.00}{#1}}
\newcommand{\RegionMarkerTok}[1]{\textcolor[rgb]{0.00,0.23,0.31}{#1}}
\newcommand{\SpecialCharTok}[1]{\textcolor[rgb]{0.37,0.37,0.37}{#1}}
\newcommand{\SpecialStringTok}[1]{\textcolor[rgb]{0.13,0.47,0.30}{#1}}
\newcommand{\StringTok}[1]{\textcolor[rgb]{0.13,0.47,0.30}{#1}}
\newcommand{\VariableTok}[1]{\textcolor[rgb]{0.07,0.07,0.07}{#1}}
\newcommand{\VerbatimStringTok}[1]{\textcolor[rgb]{0.13,0.47,0.30}{#1}}
\newcommand{\WarningTok}[1]{\textcolor[rgb]{0.37,0.37,0.37}{\textit{#1}}}

\usepackage{longtable,booktabs,array}
\newcounter{none} % for unnumbered tables
\usepackage{calc} % for calculating minipage widths
% Correct order of tables after \paragraph or \subparagraph
\usepackage{etoolbox}
\makeatletter
\patchcmd\longtable{\par}{\if@noskipsec\mbox{}\fi\par}{}{}
\makeatother
% Allow footnotes in longtable head/foot
\IfFileExists{footnotehyper.sty}{\usepackage{footnotehyper}}{\usepackage{footnote}}
\makesavenoteenv{longtable}
\usepackage{graphicx}
\makeatletter
\newsavebox\pandoc@box
\newcommand*\pandocbounded[1]{% scales image to fit in text height/width
  \sbox\pandoc@box{#1}%
  \Gscale@div\@tempa{\textheight}{\dimexpr\ht\pandoc@box+\dp\pandoc@box\relax}%
  \Gscale@div\@tempb{\linewidth}{\wd\pandoc@box}%
  \ifdim\@tempb\p@<\@tempa\p@\let\@tempa\@tempb\fi% select the smaller of both
  \ifdim\@tempa\p@<\p@\scalebox{\@tempa}{\usebox\pandoc@box}%
  \else\usebox{\pandoc@box}%
  \fi%
}
% Set default figure placement to htbp
\def\fps@figure{htbp}
\makeatother





\setlength{\emergencystretch}{3em} % prevent overfull lines

\providecommand{\tightlist}{%
  \setlength{\itemsep}{0pt}\setlength{\parskip}{0pt}}



 


\KOMAoption{captions}{tableheading}
\makeatletter
\@ifpackageloaded{caption}{}{\usepackage{caption}}
\AtBeginDocument{%
\ifdefined\contentsname
  \renewcommand*\contentsname{Table of contents}
\else
  \newcommand\contentsname{Table of contents}
\fi
\ifdefined\listfigurename
  \renewcommand*\listfigurename{List of Figures}
\else
  \newcommand\listfigurename{List of Figures}
\fi
\ifdefined\listtablename
  \renewcommand*\listtablename{List of Tables}
\else
  \newcommand\listtablename{List of Tables}
\fi
\ifdefined\figurename
  \renewcommand*\figurename{Figure}
\else
  \newcommand\figurename{Figure}
\fi
\ifdefined\tablename
  \renewcommand*\tablename{Table}
\else
  \newcommand\tablename{Table}
\fi
}
\@ifpackageloaded{float}{}{\usepackage{float}}
\floatstyle{ruled}
\@ifundefined{c@chapter}{\newfloat{codelisting}{h}{lop}}{\newfloat{codelisting}{h}{lop}[chapter]}
\floatname{codelisting}{Listing}
\newcommand*\listoflistings{\listof{codelisting}{List of Listings}}
\makeatother
\makeatletter
\makeatother
\makeatletter
\@ifpackageloaded{caption}{}{\usepackage{caption}}
\@ifpackageloaded{subcaption}{}{\usepackage{subcaption}}
\makeatother
\usepackage{bookmark}
\IfFileExists{xurl.sty}{\usepackage{xurl}}{} % add URL line breaks if available
\urlstyle{same}
% fallback for those not using the hyperref driver hyperxmp:
\makeatletter
\@ifundefined{xmpquote}{\newcommand{\xmpquote}[1]{#1}}{}
\makeatother
\hypersetup{
  pdftitle={Quality control --- Study B (contraception)},
  colorlinks=true,
  linkcolor={blue},
  filecolor={Maroon},
  citecolor={Blue},
  urlcolor={Blue},
  pdfcreator={LaTeX via pandoc}}


\title{Quality control --- Study B (contraception)}
\author{}
\date{}
\begin{document}
\maketitle


\begin{quote}
\textbf{What this document is.} Before we test anything, we look hard at
the raw data and ask: are the three files consistent with each other,
are there samples that behave strangely, and is anything broken badly
enough that we should drop or fix it now? Every number below was
computed from the data --- at the bottom of each section there is a
\emph{sanity check} telling you how to spot-check it yourself.
\end{quote}

\begin{Shaded}
\begin{Highlighting}[]
\FunctionTok{library}\NormalTok{(tidyverse)}
\FunctionTok{library}\NormalTok{(patchwork)   }\CommentTok{\# puts several ggplots in one row}

\CommentTok{\# Workshop helper: read\_counts() loads the three count tables plus metadata}
\CommentTok{\# and cytokines, and stops with an error if the sample ids do not line up.}
\FunctionTok{source}\NormalTok{(here}\SpecialCharTok{::}\FunctionTok{here}\NormalTok{(}\StringTok{"R/count\_tables.R"}\NormalTok{))}

\NormalTok{d  }\OtherTok{\textless{}{-}} \FunctionTok{read\_counts}\NormalTok{(}\StringTok{"group\_b"}\NormalTok{)      }\CommentTok{\# the contraception study}
\NormalTok{sp }\OtherTok{\textless{}{-}}\NormalTok{ d}\SpecialCharTok{$}\NormalTok{species                   }\CommentTok{\# 220 species  x 93 samples}
\NormalTok{ge }\OtherTok{\textless{}{-}}\NormalTok{ d}\SpecialCharTok{$}\NormalTok{genes                     }\CommentTok{\# 14,226 genes x 93 samples}
\NormalTok{md }\OtherTok{\textless{}{-}}\NormalTok{ d}\SpecialCharTok{$}\NormalTok{metadata                  }\CommentTok{\# one row per sample}
\NormalTok{cy }\OtherTok{\textless{}{-}}\NormalTok{ d}\SpecialCharTok{$}\NormalTok{cytokines                 }\CommentTok{\# 15 immune markers, one row per sample}
\end{Highlighting}
\end{Shaded}

\section{1. Do the pieces fit
together?}\label{do-the-pieces-fit-together}

The most dangerous QC failure is silent mis-alignment: if the
count-table columns and the metadata rows were in a different order,
every test would be wrong and nothing would warn us.
\texttt{read\_counts()} refuses to return if the sample ids of the two
count tables and the metadata are not exactly the same set, and it
reorders the columns to match the metadata --- but it does \textbf{not}
check the cytokine table, so we check all four tables ourselves here,
both for the same \emph{set} of ids and for the same \emph{order}.

\begin{Shaded}
\begin{Highlighting}[]
\CommentTok{\# Every cell is turned into text, so one column can hold counts and yes/no answers.}
\FunctionTok{tibble}\NormalTok{(}
  \AttributeTok{check =} \FunctionTok{c}\NormalTok{(}
    \StringTok{"samples in metadata"}\NormalTok{,}
    \StringTok{"species table: features x samples"}\NormalTok{,}
    \StringTok{"gene table: features x samples"}\NormalTok{,}
    \StringTok{"cytokines rows"}\NormalTok{,}
    \StringTok{"sample\_id appears more than once?"}\NormalTok{,}
    \StringTok{"participant appears more than once?"}\NormalTok{,}
    \StringTok{"species ids same set as metadata?"}\NormalTok{,}
    \StringTok{"gene ids same set as metadata?"}\NormalTok{,}
    \StringTok{"cytokine ids same set as metadata?"}\NormalTok{,}
    \StringTok{"species columns ordered like metadata?"}\NormalTok{,}
    \StringTok{"gene columns ordered like metadata?"}\NormalTok{,}
    \StringTok{"cytokine rows ordered like metadata?"}\NormalTok{,}
    \StringTok{"negative or fractional counts anywhere?"}\NormalTok{,}
    \StringTok{"non{-}integer counts?"}
\NormalTok{  ),}
  \AttributeTok{result =} \FunctionTok{c}\NormalTok{(}
    \FunctionTok{nrow}\NormalTok{(md),}
    \FunctionTok{str\_c}\NormalTok{(}\FunctionTok{nrow}\NormalTok{(sp), }\StringTok{" x "}\NormalTok{, }\FunctionTok{ncol}\NormalTok{(sp)),}
    \FunctionTok{str\_c}\NormalTok{(}\FunctionTok{nrow}\NormalTok{(ge), }\StringTok{" x "}\NormalTok{, }\FunctionTok{ncol}\NormalTok{(ge)),}
    \FunctionTok{nrow}\NormalTok{(cy),}
    \FunctionTok{anyDuplicated}\NormalTok{(md}\SpecialCharTok{$}\NormalTok{sample\_id),}
    \FunctionTok{anyDuplicated}\NormalTok{(md}\SpecialCharTok{$}\NormalTok{participant\_id),}
    \FunctionTok{setequal}\NormalTok{(md}\SpecialCharTok{$}\NormalTok{sample\_id, }\FunctionTok{colnames}\NormalTok{(sp)),}
    \FunctionTok{setequal}\NormalTok{(md}\SpecialCharTok{$}\NormalTok{sample\_id, }\FunctionTok{colnames}\NormalTok{(ge)),}
    \FunctionTok{setequal}\NormalTok{(md}\SpecialCharTok{$}\NormalTok{sample\_id, cy}\SpecialCharTok{$}\NormalTok{sample\_id),}
    \FunctionTok{identical}\NormalTok{(}\FunctionTok{colnames}\NormalTok{(sp), md}\SpecialCharTok{$}\NormalTok{sample\_id),}
    \FunctionTok{identical}\NormalTok{(}\FunctionTok{colnames}\NormalTok{(ge), md}\SpecialCharTok{$}\NormalTok{sample\_id),}
    \FunctionTok{identical}\NormalTok{(cy}\SpecialCharTok{$}\NormalTok{sample\_id, md}\SpecialCharTok{$}\NormalTok{sample\_id),}
    \FunctionTok{any}\NormalTok{(sp }\SpecialCharTok{\textless{}} \DecValTok{0}\NormalTok{) }\SpecialCharTok{||} \FunctionTok{any}\NormalTok{(ge }\SpecialCharTok{\textless{}} \DecValTok{0}\NormalTok{),}
    \FunctionTok{any}\NormalTok{(sp }\SpecialCharTok{!=} \FunctionTok{round}\NormalTok{(sp)) }\SpecialCharTok{||} \FunctionTok{any}\NormalTok{(ge }\SpecialCharTok{!=} \FunctionTok{round}\NormalTok{(ge))}
\NormalTok{  ) }\SpecialCharTok{|\textgreater{}}
    \FunctionTok{as.character}\NormalTok{()}
\NormalTok{)}
\end{Highlighting}
\end{Shaded}

\begin{verbatim}
# A tibble: 14 x 2
   check                                   result    
   <chr>                                   <chr>     
 1 samples in metadata                     93        
 2 species table: features x samples       220 x 93  
 3 gene table: features x samples          14226 x 93
 4 cytokines rows                          93        
 5 sample_id appears more than once?       0         
 6 participant appears more than once?     0         
 7 species ids same set as metadata?       TRUE      
 8 gene ids same set as metadata?          TRUE      
 9 cytokine ids same set as metadata?      TRUE      
10 species columns ordered like metadata?  TRUE      
11 gene columns ordered like metadata?     TRUE      
12 cytokine rows ordered like metadata?    TRUE      
13 negative or fractional counts anywhere? FALSE     
14 non-integer counts?                     FALSE     
\end{verbatim}

\begin{quote}
\textbf{Sanity check.} The workshop materials say Study B has 93
samples. Both tables should say \texttt{x\ 93}. One participant should
give one sample --- the duplicates check must be \texttt{0}, and all six
id checks should say \texttt{TRUE}. To verify one of them yourself: open
\texttt{06\_cytokines.csv}, sort its first column, and confirm the first
and last entries are \texttt{B001} and \texttt{B093}, just as in
\texttt{01\_metadata.csv}.
\end{quote}

\section{2. Who is in the study, and is it
balanced?}\label{who-is-in-the-study-and-is-it-balanced}

Downstream tests assume the groups we compare are big enough and not
wildly lopsided. Also worth checking now: is \texttt{rnaseq\_batch}
(which sequencing run a sample was processed in) mixed evenly across the
arms, or confounded with them? If batch sat entirely inside one arm we
could never separate ``the pill changed the genes'' from ``the second
run did''.

\begin{Shaded}
\begin{Highlighting}[]
\NormalTok{knitr}\SpecialCharTok{::}\FunctionTok{kable}\NormalTok{(}\FunctionTok{table}\NormalTok{(}\AttributeTok{Arm =}\NormalTok{ md}\SpecialCharTok{$}\NormalTok{arm, }\StringTok{\textasciigrave{}}\AttributeTok{RNA{-}seq batch}\StringTok{\textasciigrave{}} \OtherTok{=}\NormalTok{ md}\SpecialCharTok{$}\NormalTok{rnaseq\_batch))}
\end{Highlighting}
\end{Shaded}

{\def\LTcaptype{none} % do not increment counter
\begin{longtable}[]{@{}lrr@{}}
\toprule\noalign{}
& B1 & B2 \\
\midrule\noalign{}
\endhead
\bottomrule\noalign{}
\endlastfoot
injectable & 20 & 19 \\
pill & 14 & 14 \\
ring & 13 & 13 \\
\end{longtable}
}

All 93 participants contribute exactly one sample. The three arms (93
samples: 28 pill, 39 injectable, 26 ring) are reasonably close in size
and both sequencing batches contain samples from all three arms, so
batch can be included as a covariate without swallowing the treatment.

Study B has a single visit per participant, so there is no time point to
cross-tabulate against \texttt{rnaseq\_batch} --- that check belongs to
Study A. What we can ask here is the Study B version of ``did the arms
differ at baseline?'': with one visit, every metadata column is a
baseline measurement, so we compare them across arms.

\begin{Shaded}
\begin{Highlighting}[]
\NormalTok{knitr}\SpecialCharTok{::}\FunctionTok{kable}\NormalTok{(}\FunctionTok{table}\NormalTok{(}\AttributeTok{Arm =}\NormalTok{ md}\SpecialCharTok{$}\NormalTok{arm, }\StringTok{\textasciigrave{}}\AttributeTok{BV status}\StringTok{\textasciigrave{}} \OtherTok{=}\NormalTok{ md}\SpecialCharTok{$}\NormalTok{bv))}
\end{Highlighting}
\end{Shaded}

{\def\LTcaptype{none} % do not increment counter
\begin{longtable}[]{@{}lrrr@{}}
\toprule\noalign{}
& intermediate & negative & positive \\
\midrule\noalign{}
\endhead
\bottomrule\noalign{}
\endlastfoot
injectable & 0 & 20 & 19 \\
pill & 2 & 16 & 10 \\
ring & 1 & 15 & 10 \\
\end{longtable}
}

\begin{Shaded}
\begin{Highlighting}[]
\NormalTok{knitr}\SpecialCharTok{::}\FunctionTok{kable}\NormalTok{(}\FunctionTok{table}\NormalTok{(}\AttributeTok{Arm =}\NormalTok{ md}\SpecialCharTok{$}\NormalTok{arm, }\StringTok{\textasciigrave{}}\AttributeTok{Any STI}\StringTok{\textasciigrave{}} \OtherTok{=}\NormalTok{ md}\SpecialCharTok{$}\NormalTok{any\_sti))}
\end{Highlighting}
\end{Shaded}

{\def\LTcaptype{none} % do not increment counter
\begin{longtable}[]{@{}lrr@{}}
\toprule\noalign{}
& 0 & 1 \\
\midrule\noalign{}
\endhead
\bottomrule\noalign{}
\endlastfoot
injectable & 20 & 19 \\
pill & 16 & 12 \\
ring & 9 & 17 \\
\end{longtable}
}

\begin{Shaded}
\begin{Highlighting}[]
\FunctionTok{tibble}\NormalTok{(}
  \AttributeTok{variable =} \FunctionTok{c}\NormalTok{(}\StringTok{"Nugent score"}\NormalTok{, }\StringTok{"pH"}\NormalTok{),}
  \AttributeTok{test     =} \FunctionTok{c}\NormalTok{(}\StringTok{"Kruskal–Wallis"}\NormalTok{, }\StringTok{"Kruskal–Wallis"}\NormalTok{),}
  \AttributeTok{p\_value  =} \FunctionTok{signif}\NormalTok{(}\FunctionTok{c}\NormalTok{(}\FunctionTok{kruskal.test}\NormalTok{(nugent }\SpecialCharTok{\textasciitilde{}}\NormalTok{ arm, md)}\SpecialCharTok{$}\NormalTok{p.value,}
                      \FunctionTok{kruskal.test}\NormalTok{(ph }\SpecialCharTok{\textasciitilde{}}\NormalTok{ arm, md)}\SpecialCharTok{$}\NormalTok{p.value), }\DecValTok{2}\NormalTok{)}
\NormalTok{) }\SpecialCharTok{|\textgreater{}}\NormalTok{ knitr}\SpecialCharTok{::}\FunctionTok{kable}\NormalTok{()}
\end{Highlighting}
\end{Shaded}

{\def\LTcaptype{none} % do not increment counter
\begin{longtable}[]{@{}llr@{}}
\toprule\noalign{}
variable & test & p\_value \\
\midrule\noalign{}
\endhead
\bottomrule\noalign{}
\endlastfoot
Nugent score & Kruskal--Wallis & 0.63 \\
pH & Kruskal--Wallis & 0.02 \\
\end{longtable}
}

BV status (Fisher's exact p = 0.45) and STI history (p = 0.22) are
distributed similarly across arms, and Nugent scores are too (median 2
in pill and injectable, 2.5 in ring). \textbf{Vaginal pH does differ}
(Kruskal--Wallis p = 0.02; median 4.55 in the pill arm vs 5.0 in
injectable and ring, and the ring arm has no sample below 4.4). The
shift is small and could easily be a chance imbalance among 26--39
people, but because pH tracks BV and shifts with contraceptive hormones
in the literature, we will keep it in mind --- and consider adjusting
for it --- when the arms are compared on the microbiome.

\section{3. Missing values}\label{missing-values}

Count tables almost never have missing values (an unmeasured gene just
gets zero reads), so missingness lives in the spreadsheet-like tables.
For the metadata the answer is good news: \textbf{not a single cell is
missing} --- every one of the 8 columns (\texttt{sample\_id},
\texttt{participant\_id}, \texttt{arm}, \texttt{bv}, \texttt{nugent},
\texttt{ph}, \texttt{any\_sti}, \texttt{rnaseq\_batch}) is filled for
all 93 samples, so no sample has to be dropped for lack of arm or batch
information and no variable needs imputation or a ``missing'' category:

\begin{Shaded}
\begin{Highlighting}[]
\FunctionTok{bind\_rows}\NormalTok{(}
  \AttributeTok{metadata  =} \FunctionTok{tibble}\NormalTok{(}\AttributeTok{source =} \StringTok{"metadata"}\NormalTok{,  }\AttributeTok{missing =} \FunctionTok{sum}\NormalTok{(}\FunctionTok{is.na}\NormalTok{(md)),  }\AttributeTok{cells =} \FunctionTok{nrow}\NormalTok{(md) }\SpecialCharTok{*} \FunctionTok{ncol}\NormalTok{(md)),}
  \AttributeTok{cytokines =} \FunctionTok{tibble}\NormalTok{(}\AttributeTok{source =} \StringTok{"cytokines"}\NormalTok{, }\AttributeTok{missing =} \FunctionTok{sum}\NormalTok{(}\FunctionTok{is.na}\NormalTok{(cy)), }\AttributeTok{cells =} \FunctionTok{nrow}\NormalTok{(cy) }\SpecialCharTok{*} \FunctionTok{ncol}\NormalTok{(cy))}
\NormalTok{) }\SpecialCharTok{|\textgreater{}}
  \FunctionTok{mutate}\NormalTok{(}\AttributeTok{percent =} \FunctionTok{round}\NormalTok{(}\DecValTok{100} \SpecialCharTok{*}\NormalTok{ missing }\SpecialCharTok{/}\NormalTok{ cells, }\DecValTok{2}\NormalTok{))}
\end{Highlighting}
\end{Shaded}

\begin{verbatim}
# A tibble: 2 x 4
  source    missing cells percent
  <chr>       <int> <int>   <dbl>
1 metadata        0   744    0   
2 cytokines      30  1488    2.02
\end{verbatim}

So the only missing data in the study is here: 30 of 1,395 cytokine
cells. Which samples?

\begin{Shaded}
\begin{Highlighting}[]
\NormalTok{md\_missing }\OtherTok{\textless{}{-}}\NormalTok{ md}\SpecialCharTok{$}\NormalTok{sample\_id[}\FunctionTok{rowSums}\NormalTok{(}\FunctionTok{is.na}\NormalTok{(cy[, }\SpecialCharTok{{-}}\DecValTok{1}\NormalTok{])) }\SpecialCharTok{\textgreater{}} \DecValTok{0}\NormalTok{]}
\NormalTok{md\_missing}
\end{Highlighting}
\end{Shaded}

\begin{verbatim}
[1] "B018" "B029"
\end{verbatim}

\textbf{Finding: \texttt{B018} and \texttt{B029} are missing all 15
cytokines} (the columns are skipped, not partly filled), which reads
like an assay that failed for two samples. Nothing to impute --- they
simply have no immune measurements. We keep them for the microbiome and
gene analyses and lose them only in cytokine work.

\begin{quote}
\textbf{Sanity check.} \texttt{is.na(cy)} should count exactly 30 (2
samples × 15 markers), and the two ids above should appear in
\texttt{01\_metadata.csv} with all other fields filled in.
\end{quote}

\section{4. Library sizes: how much did we
sequence?}\label{library-sizes-how-much-did-we-sequence}

``Library size'' means the total number of reads mapped for a sample ---
literally the column sum of the count table. Samples far below the rest
may just have been sequenced shallowly, which makes them noisy, not
broken.

\begin{Shaded}
\begin{Highlighting}[]
\NormalTok{qc }\OtherTok{\textless{}{-}} \FunctionTok{bind\_cols}\NormalTok{(}
\NormalTok{  md,}
  \AttributeTok{lib\_species =} \FunctionTok{colSums}\NormalTok{(sp),}
  \AttributeTok{lib\_genes   =} \FunctionTok{colSums}\NormalTok{(ge)}
\NormalTok{)}

\NormalTok{p\_sp }\OtherTok{\textless{}{-}} \FunctionTok{ggplot}\NormalTok{(qc, }\FunctionTok{aes}\NormalTok{(lib\_species)) }\SpecialCharTok{+}
  \FunctionTok{geom\_histogram}\NormalTok{(}\AttributeTok{bins =} \DecValTok{30}\NormalTok{, }\AttributeTok{fill =} \StringTok{"grey70"}\NormalTok{, }\AttributeTok{colour =} \StringTok{"white"}\NormalTok{) }\SpecialCharTok{+}
  \FunctionTok{labs}\NormalTok{(}\AttributeTok{title =} \StringTok{"Microbiome"}\NormalTok{, }\AttributeTok{x =} \StringTok{"reads"}\NormalTok{, }\AttributeTok{y =} \StringTok{"samples"}\NormalTok{)}

\NormalTok{p\_ge }\OtherTok{\textless{}{-}} \FunctionTok{ggplot}\NormalTok{(qc, }\FunctionTok{aes}\NormalTok{(lib\_genes)) }\SpecialCharTok{+}
  \FunctionTok{geom\_histogram}\NormalTok{(}\AttributeTok{bins =} \DecValTok{30}\NormalTok{, }\AttributeTok{fill =} \StringTok{"grey70"}\NormalTok{, }\AttributeTok{colour =} \StringTok{"white"}\NormalTok{) }\SpecialCharTok{+}
  \FunctionTok{labs}\NormalTok{(}\AttributeTok{title =} \StringTok{"RNA{-}seq"}\NormalTok{, }\AttributeTok{x =} \StringTok{"reads"}\NormalTok{, }\AttributeTok{y =} \ConstantTok{NULL}\NormalTok{)}

\NormalTok{p\_sp }\SpecialCharTok{+}\NormalTok{ p\_ge}
\end{Highlighting}
\end{Shaded}

\pandocbounded{\includegraphics[keepaspectratio]{01-qc_files/figure-pdf/libsize-1.pdf}}

\begin{Shaded}
\begin{Highlighting}[]
\CommentTok{\# The rules below are our choices, not universal constants — see "Choices made".}
\NormalTok{rule }\OtherTok{\textless{}{-}} \FunctionTok{median}\NormalTok{(qc}\SpecialCharTok{$}\NormalTok{lib\_species) }\SpecialCharTok{/} \DecValTok{2}
\NormalTok{qc }\SpecialCharTok{|\textgreater{}} \FunctionTok{filter}\NormalTok{(lib\_species }\SpecialCharTok{\textless{}}\NormalTok{ rule) }\SpecialCharTok{|\textgreater{}} \FunctionTok{select}\NormalTok{(sample\_id, lib\_species) }\SpecialCharTok{|\textgreater{}}
\NormalTok{  knitr}\SpecialCharTok{::}\FunctionTok{kable}\NormalTok{(}\AttributeTok{digits =} \DecValTok{0}\NormalTok{, }\AttributeTok{caption =} \FunctionTok{str\_c}\NormalTok{(}\StringTok{"Microbiome: samples below half the median library size ("}\NormalTok{, }\FunctionTok{format}\NormalTok{(rule, }\AttributeTok{big.mark =} \StringTok{","}\NormalTok{), }\StringTok{" reads)"}\NormalTok{))}
\end{Highlighting}
\end{Shaded}

\begin{longtable}[]{@{}lr@{}}
\caption{Microbiome: samples below half the median library size (14,606
reads)}\tabularnewline
\toprule\noalign{}
sample\_id & lib\_species \\
\midrule\noalign{}
\endfirsthead
\toprule\noalign{}
sample\_id & lib\_species \\
\midrule\noalign{}
\endhead
\bottomrule\noalign{}
\endlastfoot
B003 & 14057 \\
B006 & 5247 \\
B024 & 11627 \\
B065 & 11570 \\
B069 & 12927 \\
B070 & 13333 \\
B076 & 14534 \\
B083 & 10756 \\
B091 & 13543 \\
\end{longtable}

Species libraries run from about 5,000 to 77,000 reads (median 29,000),
and nine samples fall below half the median --- \texttt{B006} is the
lowest at 5,247 reads, roughly six times fewer reads than the median
sample. None is so low that it must be dropped; we flag them and will
look again after ordination in case any sits alone. Gene libraries are
around 6--24 million reads, tightly packed --- only one sample sits
below half the median there, \textbf{\texttt{B042} at 5.9 million}
reads, and only barely (the rule is 6.4 million). \texttt{B042} is not
one of the nine low-depth microbiome samples, so its shallow sequencing
is specific to the RNA-seq library, not to the sample as a whole.

\begin{quote}
\textbf{Sanity check.} Add the first five columns of any sample of your
choice in \texttt{02\_microbiome\_counts.csv} with a calculator --- it
should match \texttt{qc\$lib\_species{[}qc\$sample\_id\ ==\ "B001"{]}}
(or whichever sample you picked).
\end{quote}

\section{5. Dominance: is one feature eating each
sample?}\label{dominance-is-one-feature-eating-each-sample}

Both tables are \textbf{compositional}: sequencing gives us proportions,
not absolute amounts, so every table must sum to the library size. That
means a huge read count for one feature mechanically shrinks all the
others. Some dominance is expected --- vaginal communities are often 90
\% one \emph{Lactobacillus} --- but near-total domination is worth
seeing explicitly.

\begin{Shaded}
\begin{Highlighting}[]
\NormalTok{top\_share }\OtherTok{\textless{}{-}} \ControlFlowTok{function}\NormalTok{(m) }\FunctionTok{apply}\NormalTok{(m, }\DecValTok{2}\NormalTok{, max) }\SpecialCharTok{/} \FunctionTok{colSums}\NormalTok{(m)}
\NormalTok{qc}\SpecialCharTok{$}\NormalTok{top\_species }\OtherTok{\textless{}{-}} \FunctionTok{top\_share}\NormalTok{(sp)}
\NormalTok{qc}\SpecialCharTok{$}\NormalTok{top\_gene    }\OtherTok{\textless{}{-}} \FunctionTok{top\_share}\NormalTok{(ge)}

\NormalTok{worst\_sp }\OtherTok{\textless{}{-}}\NormalTok{ qc }\SpecialCharTok{|\textgreater{}} \FunctionTok{slice\_max}\NormalTok{(top\_species, }\AttributeTok{n =} \DecValTok{3}\NormalTok{) }\SpecialCharTok{|\textgreater{}}
  \FunctionTok{transmute}\NormalTok{(sample\_id, }\AttributeTok{share =} \FunctionTok{str\_c}\NormalTok{(}\FunctionTok{round}\NormalTok{(}\DecValTok{100} \SpecialCharTok{*}\NormalTok{ top\_species, }\DecValTok{1}\NormalTok{), }\StringTok{"\%"}\NormalTok{))}

\NormalTok{knitr}\SpecialCharTok{::}\FunctionTok{kable}\NormalTok{(worst\_sp, }\AttributeTok{caption =} \StringTok{"Samples whose single most abundant species takes the largest share"}\NormalTok{)}
\end{Highlighting}
\end{Shaded}

\begin{longtable}[]{@{}ll@{}}
\caption{Samples whose single most abundant species takes the largest
share}\tabularnewline
\toprule\noalign{}
sample\_id & share \\
\midrule\noalign{}
\endfirsthead
\toprule\noalign{}
sample\_id & share \\
\midrule\noalign{}
\endhead
\bottomrule\noalign{}
\endlastfoot
B030 & 99.8\% \\
B016 & 99.7\% \\
B012 & 99.7\% \\
\end{longtable}

The median sample's biggest species is about half of all reads ---
normal for this body site. The extreme is \texttt{B030}, which is 99.8
\% \emph{Lactobacillus iners}. That is not an error: some women really
are monoculture, and it is exactly the biology we want to study. We keep
it, but we will not be surprised later if it pulls ordinations towards
itself.

\begin{quote}
\textbf{Sanity check.} In \texttt{02\_microbiome\_counts.csv}, find the
\texttt{B030} column, sort it, and confirm the biggest number is about
99.8 \% of the column sum.
\end{quote}

\section{6. Anything odd inside the
RNA-seq?}\label{anything-odd-inside-the-rna-seq}

The largest gene share in any sample is 29\% --- in \texttt{B093},
nearly a third of all reads come from \textbf{HBB}, the beta chain of
haemoglobin. Haemoglobin genes are the classic signature of blood
contamination in mucosal swabs: the sample picked up red blood cells,
and RNA from blood crowded out cervical/vaginal transcripts. It is worth
remembering when interpreting \texttt{B093} specifically, but a quarter
of samples carry some HBB signal and the difference is quantitative, not
broken, so we keep the sample.

\begin{Shaded}
\begin{Highlighting}[]
\CommentTok{\# Proportion of each sample\textquotesingle{}s reads belonging to the haemoglobin genes HBB/HBA}
\NormalTok{hb }\OtherTok{\textless{}{-}}\NormalTok{ d}\SpecialCharTok{$}\NormalTok{gene\_info }\SpecialCharTok{|\textgreater{}} \FunctionTok{filter}\NormalTok{(symbol }\SpecialCharTok{\%in\%} \FunctionTok{c}\NormalTok{(}\StringTok{"HBB"}\NormalTok{, }\StringTok{"HBA1"}\NormalTok{, }\StringTok{"HBA2"}\NormalTok{)) }\SpecialCharTok{|\textgreater{}} \FunctionTok{pull}\NormalTok{(ensembl\_gene\_id)}
\NormalTok{qc}\SpecialCharTok{$}\NormalTok{hbb\_share }\OtherTok{\textless{}{-}} \FunctionTok{colSums}\NormalTok{(ge[hb, , }\AttributeTok{drop =} \ConstantTok{FALSE}\NormalTok{]) }\SpecialCharTok{/}\NormalTok{ qc}\SpecialCharTok{$}\NormalTok{lib\_genes}

\FunctionTok{ggplot}\NormalTok{(qc, }\FunctionTok{aes}\NormalTok{(}\FunctionTok{reorder}\NormalTok{(sample\_id, hbb\_share), hbb\_share)) }\SpecialCharTok{+}
  \FunctionTok{geom\_col}\NormalTok{(}\AttributeTok{fill =} \StringTok{"grey70"}\NormalTok{) }\SpecialCharTok{+}
  \FunctionTok{scale\_y\_continuous}\NormalTok{(}\AttributeTok{labels =}\NormalTok{ scales}\SpecialCharTok{::}\NormalTok{percent) }\SpecialCharTok{+}
  \FunctionTok{labs}\NormalTok{(}\AttributeTok{title =} \StringTok{"Haemoglobin (HBB/HBA) reads per sample"}\NormalTok{,}
       \AttributeTok{x =} \ConstantTok{NULL}\NormalTok{, }\AttributeTok{y =} \StringTok{"share of all gene reads"}\NormalTok{) }\SpecialCharTok{+}
  \FunctionTok{theme}\NormalTok{(}\AttributeTok{axis.text.x =} \FunctionTok{element\_blank}\NormalTok{())}
\end{Highlighting}
\end{Shaded}

\pandocbounded{\includegraphics[keepaspectratio]{01-qc_files/figure-pdf/hbb-1.pdf}}

\begin{quote}
\textbf{Sanity check.} Sort one sample's column in
\texttt{04\_rnaseq\_counts.csv.gz} (unzip it first) by value; the top
gene for \texttt{B093} should be HBB.
\end{quote}

\section{7. Batch effects at a glance}\label{batch-effects-at-a-glance}

\begin{Shaded}
\begin{Highlighting}[]
\FunctionTok{ggplot}\NormalTok{(qc, }\FunctionTok{aes}\NormalTok{(rnaseq\_batch, lib\_genes, }\AttributeTok{fill =}\NormalTok{ arm)) }\SpecialCharTok{+}
  \FunctionTok{geom\_boxplot}\NormalTok{() }\SpecialCharTok{+}
  \FunctionTok{scale\_y\_log10}\NormalTok{(}\AttributeTok{labels =}\NormalTok{ scales}\SpecialCharTok{::}\NormalTok{comma) }\SpecialCharTok{+}
  \FunctionTok{labs}\NormalTok{(}\AttributeTok{title =} \StringTok{"Gene library size by sequencing batch and arm"}\NormalTok{,}
       \AttributeTok{y =} \StringTok{"reads (log10)"}\NormalTok{, }\AttributeTok{fill =} \StringTok{"Arm"}\NormalTok{, }\AttributeTok{x =} \StringTok{"Sequencing batch"}\NormalTok{)}
\end{Highlighting}
\end{Shaded}

\pandocbounded{\includegraphics[keepaspectratio]{01-qc_files/figure-pdf/batch-1.pdf}}

No batch jumps out as different in depth, and both arms appear in both
batches. A real check for batch effects in expression profiles comes
later (ordination coloured by batch); this only rules out crude depth
differences.

\section{8. How prevalent are the
features?}\label{how-prevalent-are-the-features}

\begin{Shaded}
\begin{Highlighting}[]
\NormalTok{prev\_sp }\OtherTok{\textless{}{-}} \FunctionTok{rowSums}\NormalTok{(sp }\SpecialCharTok{\textgreater{}} \DecValTok{0}\NormalTok{) }\SpecialCharTok{/} \FunctionTok{ncol}\NormalTok{(sp)}
\NormalTok{prev\_ge }\OtherTok{\textless{}{-}} \FunctionTok{rowSums}\NormalTok{(ge }\SpecialCharTok{\textgreater{}} \DecValTok{0}\NormalTok{) }\SpecialCharTok{/} \FunctionTok{ncol}\NormalTok{(ge)}

\FunctionTok{data.frame}\NormalTok{(}\AttributeTok{prevalence =} \FunctionTok{c}\NormalTok{(prev\_sp, prev\_ge),}
           \AttributeTok{table =} \FunctionTok{rep}\NormalTok{(}\FunctionTok{c}\NormalTok{(}\StringTok{"species"}\NormalTok{, }\StringTok{"genes"}\NormalTok{), }\FunctionTok{c}\NormalTok{(}\FunctionTok{length}\NormalTok{(prev\_sp), }\FunctionTok{length}\NormalTok{(prev\_ge)))) }\SpecialCharTok{|\textgreater{}}
  \FunctionTok{ggplot}\NormalTok{(}\FunctionTok{aes}\NormalTok{(prevalence, }\AttributeTok{fill =}\NormalTok{ table)) }\SpecialCharTok{+}
  \FunctionTok{geom\_density}\NormalTok{(}\AttributeTok{alpha =} \FloatTok{0.5}\NormalTok{) }\SpecialCharTok{+}
  \FunctionTok{labs}\NormalTok{(}\AttributeTok{title =} \StringTok{"In how many samples is each feature seen?"}\NormalTok{,}
       \AttributeTok{x =} \StringTok{"fraction of samples with the feature present"}\NormalTok{, }\AttributeTok{y =} \ConstantTok{NULL}\NormalTok{, }\AttributeTok{fill =} \ConstantTok{NULL}\NormalTok{)}
\end{Highlighting}
\end{Shaded}

\pandocbounded{\includegraphics[keepaspectratio]{01-qc_files/figure-pdf/prevalence-1.pdf}}

Most species are rare (72 \% of them appear in under 10 \% of samples)
--- typical for microbial data, and a reminder that feature filtering
(\texttt{min\_samples}) will discard a lot before any
differential-abundance test. Genes, in contrast, are dense: nearly all
are detected in most samples, because sequencing 14 M reads into 14,226
well-known genes rarely leaves one empty.

\section{Verdict}\label{verdict}

\begin{itemize}
\tightlist
\item
  \textbf{Structure:} all four tables (metadata, species, genes,
  cytokines) agree on the same 93 samples and 93 participants, in the
  same order; counts are non-negative integers. No structural problems.
\item
  \textbf{Balance:} 28 pill / 39 injectable / 26 ring, both sequencing
  batches present in every arm. Baseline variables are balanced across
  arms except a small difference in vaginal pH (p = 0.02) worth
  remembering at analysis time.
\item
  \textbf{Missing data:} the metadata is 100 \% complete. \texttt{B018}
  and \texttt{B029} lack all cytokines (keep for the rest; drop them for
  cytokine analyses). No other missing values anywhere.
\item
  \textbf{Flagged samples:} nine low-depth microbiome samples (still
  usable); \texttt{B030} is a near-monoculture of \emph{L. iners} (real
  biology, watch it in ordinations); \texttt{B093} carries heavy blood
  signal in its transcriptome (remember when interpreting it).
\item
  \textbf{Nothing so broken that it must be removed before analysis.}
\end{itemize}

\subsection{Choices made}\label{choices-made}

\begin{itemize}
\tightlist
\item
  Flagged low-depth samples at \textbf{below half the median library
  size}, rather than a fixed read count, so the rule adapts to the
  table's own spread; a fixed floor (e.g.~10,000 reads) would be more
  standard in published microbiome pipelines.
\item
  Did \textbf{not remove any samples} in the QC document itself ---
  removal decisions are deferred to the analysis documents, where we can
  see whether a flagged sample actually drives a result.
\item
  Did \textbf{not impute the 30 missing cytokine values} ---
  complete-case analysis (dropping the two affected samples only for
  cytokine questions) is simpler and safer.
\item
  Used \textbf{prevalence curves} to preview feature filtering but
  applied no filter here; the filter thresholds belong to each
  analysis's ``Choices made''.
\item
  Checked batch only via \textbf{library size}, not full-expression
  profiles --- the stronger check comes with ordination in the analysis
  notebooks.
\item
  For the baseline check we used \textbf{Kruskal--Wallis} for Nugent
  score and pH (non-normal, three groups) and \textbf{Fisher's exact
  test} for BV and STI (small counts in some cells); a chi-squared test
  would be the alternative, but some expected counts are too small for
  it to be trustworthy. We report the pH difference rather than adjust
  for it here --- the decision to include pH as a covariate belongs to
  each analysis document's ``Choices made''.
\end{itemize}




\end{document}
